\documentclass{oci}
\usepackage[utf8]{inputenc}
\usepackage{lipsum}
\usepackage{subcaption}

\title{Azulejos}

\begin{document}
\begin{problemDescription}
Varideri acaba de adquirir uno de los más respetados restaurantes en Portugal,
famoso por sus hermosos diseños de azulejos.
Paradójicamente, en su primera inspección notó que la calidad de la comida era
notablemente inferior a la de otros restaurantes por la zona, por lo que decidió
indagar en el asunto.
En la inspección, notó que el embaldosado de la cocina está formado por dos
tipos de azulejos, $0$ y $1$, los cuales se muestran a continuación.

\begin{figure}[h]
  \centering
  \begin{subfigure}{0.4\textwidth}
    \centering
    \includegraphics[scale=0.3]{a.eps}
    \subcaption{azulejo de tipo 0}
  \end{subfigure}
  \begin{subfigure}{0.4\textwidth}
    \centering
    \includegraphics[scale=0.3]{b.eps} &
    \subcaption{azulejo de tipo 1}
  \end{subfigure}
\end{figure}

El diseño de estos azulejos está pensado para producir un patrón de azulejos
alternados, es decir, cada azulejo de tipo $0$ tiene un azulejo de tipo $1$
a cada uno de sus lados y viceversa.
Sin embargo, actualmente los patrones en la cocina no respetaban esa regla,
irritando así al chef y afectando la calidad de su comida.
A continuación puedes ver a la izquierda un diseño equivocado como los que están
en la cocina y a la derecha uno correcto con azulejos alternados.

\begin{figure}[h]
\setlength{\tabcolsep}{0.5pt}
\begin{subfigure}{0.45\textwidth}
  \begin{center}
    {
      \renewcommand{\arraystretch}{0.25}
      \begin{tabular}{cccc}
        \includegraphics[scale=0.3]{a.eps} &
        \includegraphics[scale=0.3]{a.eps} &
        \includegraphics[scale=0.3]{b.eps} \\
        \includegraphics[scale=0.3]{b.eps} &
        \includegraphics[scale=0.3]{b.eps} &
        \includegraphics[scale=0.3]{a.eps} \\
        \includegraphics[scale=0.3]{a.eps} &
        \includegraphics[scale=0.3]{a.eps} &
        \includegraphics[scale=0.3]{b.eps}
      \end{tabular}
    }
  \end{center}
  \subcaption{patrón incorrecto}
\end{subfigure}
\begin{subfigure}{0.45\textwidth}
  \begin{center}
    {
      \renewcommand{\arraystretch}{0.25}
      \begin{tabular}{cccc}
        \includegraphics[scale=0.3]{a.eps} &
        \includegraphics[scale=0.3]{b.eps} &
        \includegraphics[scale=0.3]{a.eps} \\
        \includegraphics[scale=0.3]{b.eps} &
        \includegraphics[scale=0.3]{a.eps} &
        \includegraphics[scale=0.3]{b.eps} \\
        \includegraphics[scale=0.3]{a.eps} &
        \includegraphics[scale=0.3]{b.eps} &
        \includegraphics[scale=0.3]{a.eps}
      \end{tabular}
    }
  \end{center}
  \subcaption{patrón correcto con azulejos alternados}
\end{subfigure}
\end{figure}

Para remediar esta situación, Vardieri ha decidido reparar los patrones, calmando
así el espíritu riguroso del chef y restaurando la calidad de sus platos.
Sin embargo, el presupuesto que Vardieri tiene disponible es limitado y necesita saber cuál es la forma más económica de realizar esta tarea.
Para cambiar un azulejo de tipo $0$ por un de tipo $1$ necesita desembolsar $P$, y para cambiar uno de tipo $1$ por uno de tipo $0$ necesita desembolsar $Q$.


¿Puedes ayudar a Varideri a encontrar la solución más económica para restaurar correctamente los azulejos?

\end{problemDescription}

\begin{inputDescription}
La entrada tiene varias líneas.

La primera línea contiene 4 enteros, $M$, $N$, $P$ y $Q$ ($1 \le M, N \le 10^3$, $1 \le P, Q \le 10^2$) que corresponden a las dimensiones $M \times N$ de un segmento rectangular de azulejos, y a los costos de reemplazar los azulejos.

Las siguientes $M$ líneas contienen $N$ números, cuyos valores pueden ser $0$ ó
$1$ indicando un tipo de azulejo.
El $i$-ésimo número de la $j$-ésima indica qué tipo de azulejo está en la
posición $(i, j)$ del embaldosado.
\end{inputDescription}

\begin{outputDescription}
La salida debe contener una línea con un único entero, el costo mínimo necesario
para cambiar el embaldosado a uno con patrón alternado.
\end{outputDescription}

\begin{scoreDescription}
  Este problema no contiene subtareas.
  Se dará puntaje proporcional a la cantidad de casos de prueba correctos siendo
  100 el puntaje máximo.
\end{scoreDescription}

\begin{sampleDescription}
\sampleIO{sample-1}
\sampleIO{sample-2}
\sampleIO{sample-3}
\end{sampleDescription}

\end{document}
